\section{Extension to exact multiplicity five}
\label{sec:m5-extension}
\subsection{Source, identifiers, and verification scope}
The extension uses the supplied archive
\path{rank7_multiplicity5_verified.zip}. Its expanded multiplication-table
file has SHA-256
\begin{quote}\small\ttfamily
\nolinkurl{a3fdcb75fbeb090913b1e4650e5124b9e0b1d016f1a71e0368a3f38e6e520c03}.
\end{quote}
It contains $10006$ rank-seven tensors through multiplicity five,
including $5805$ of exact multiplicity five. Fresh verification checks
unit, duality, Frobenius reciprocity, associativity, nonnegative integral
coefficients, and based-isomorphism uniqueness. All $7!=5040$
permutations are unnecessary: the unit must be fixed, so $6!=720$
permutations per tensor suffice. The exhaustive audit thus examines
$7204320$ unit-fixing permutations and finds no duplicates.

The $4201$ previously published tensors occur exactly in the new input.
The input does not put them in an initial contiguous block. We therefore
retain every old identifier and decision and record an explicit transport
between the previous and new physical line numbers. For the added band,
$R7\_5\_j$ denotes the $j$th row of exact multiplicity five in the supplied
file, with $1\leq j\leq5805$.
The new duality strata contain $4654$, $1012$, $138$, and $1$ rings with
respectively zero, one, two, and three nonself-dual pairs; nine of the new
rings are noncommutative.

The producer's run metadata reports complete enumeration of every
rank-seven duality stratum with no time limit. Its optimized enumeration
sources are identified by checksums but are not bundled in this input.
The present update independently verifies the output tensors, uniqueness,
and preservation of the earlier census; it does not rerun that enumeration
or independently establish its exhaustive coverage. All classification
counts for the extension are therefore statements about the supplied
$5805$ new classes, with completeness relative to all rings in the band
conditional on the producer's enumeration claim.

\subsection{A Frobenius-group realization}
\begin{proposition}\label{prop:m5-frobenius42}
The new ring $R7\_5\_5678$ has a unitary categorification
$\operatorname{Rep}(C_7\rtimes C_6)$, where a generator of $C_6$
acts on $C_7$ by multiplication by $3$.
\end{proposition}
\begin{proof}
Write the group as
$G=\langle a,b\mid a^7=b^6=1,\ bab^{-1}=a^3\rangle$.
Its six linear characters are
$\chi_j(a^xb^y)=\zeta_6^{jy}$, $0\leq j<6$.
The action of $C_6$ is transitive on the six nontrivial characters of
$\langle a\rangle$. Inducing any such character gives an irreducible
character $\rho$ of degree six, with values $6$ at the identity,
$-1$ at nonidentity elements of $\langle a\rangle$, and $0$ elsewhere.
These characters are orthonormal, and their squared degrees sum to
$6+36=42=|G|$, so they are the complete irreducible character list.
Pointwise multiplication gives
\[
 \chi_i\chi_j=\chi_{i+j\bmod6},\qquad
 \chi_i\rho=\rho,\qquad
 \rho^2=\sum_{j=0}^{5}\chi_j+5\rho.
\]
In the model order $(\chi_0,\ldots,\chi_5,\rho)$, the permutation
$(0,1,4,6,3,2,5)$ gives the corresponding census labels.
The accompanying checker verifies equality of all $343$ structure
constants under this map. Finite-group representation categories admit
a unitary structure, proving the claim. This is a known construction;
its occurrence in the extended census is not a claim of a new tensor
category.
\end{proof}

\subsection{The noncyclic maximal near-group obstruction}
The only noncommutative ring remaining after the first arithmetic
passes is $R7\_5\_4812$. Its invertible labels are
$G=\{0,1,2,3,5,6\}$ and its remaining simple is $X=b_4$, with
\[
 X^2=\sum_{g\in G}g+5X.
\]
Every invertible fixes $X$ on both sides. The group $G$ is nonabelian:
$b_1b_3=b_6$ whereas $b_3b_1=b_5$.
Nikshych--Ostrik's maximal near-group theorem
\cite[Appendix A, Theorem A.14(ii)]{OstrikRank3} states that a
near-group category with parameter $n=|G|-1$ must have cyclic $G$.
It therefore excludes this tensor for arbitrary complex fusion categories.
The appendix supplies the general near-group reduction; pivotality is
not an extra hypothesis here. The exact verifier checks the invertible
group, the entire near-group rule, and the displayed noncommuting products.

\subsection{Integral formal-codegree refinement}
\begin{proposition}\label{prop:m5-integral-codegree}
Let $R$ be a commutative fusion ring with a positive integral dimension
character $(d_i)$, and put $D=\sum_i d_i^2$.
If $R$ admits a complex categorification, then $D/f$ is a positive
integer for every formal codegree $f$ of $R$.
In particular, every formal codegree is an integer dividing $D$.
\end{proposition}
\begin{proof}
The positive character is the Frobenius--Perron character. A hypothetical
category has integral Frobenius--Perron dimension $D$ and hence is
pseudo-unitary by \cite[Proposition 8.24]{ENO}. Give it its canonical
spherical structure. Each one-dimensional character of $R$ corresponds,
by central induction, to a simple summand of the induced unit having
dimension $D/f$; see \cite[Theorem 2.13]{OstrikRank3}.
The forgetful functor preserves these dimensions, and its image is a
nonnegative integral sum of the simples of the original category. Since
all $d_i$ are integers, $D/f$ is a positive integer. Thus $f$ is rational;
as a formal codegree it is an algebraic integer, so $f$ is an integer
dividing $D$.
\end{proof}

The checker first verifies the positive dimension character against every
product and then factors the exact Casimir characteristic polynomial.
It yields the following six obstructions; only one failed factor per
ring is needed.
\begin{center}\small
\begin{tabular}{@{}l r l@{}}
\toprule
Census ID & $D$ & Obstructing codegree or factor\\
\midrule
$R7\_5\_1390$ & $112$ & $t^2-40t+80$\\
$R7\_5\_1404$ & $240$ & $t^2-40t+80$\\
$R7\_5\_2454$ & $1428$ & $t^2-48t+144$\\
$R7\_5\_4695$ & $220$ & $12$\\
$R7\_5\_5187$ & $96$ & $64$\\
$R7\_5\_5759$ & $96$ & $64$\\
\bottomrule
\end{tabular}
\end{center}
The quadratic factors are irreducible over $\mathbb Q$, whereas the
listed integer codegrees do not divide $D$. These tests strengthen the
basic $d$-number test for this integral subfamily; they do not assert
that arbitrary fusion categories have integral dimensions.

\subsection{A cyclic-stabilizer refinement}
The prime-order criterion in Proposition~\ref{prop:cyclic-module-divisibility}
admits the following extension to arbitrary cyclic order.
\begin{proposition}\label{prop:m5-cyclic-gcd}
Let $H\cong C_n$ be a subgroup of invertible objects in a complex fusion
category. Suppose that a simple $X$ is fixed on the left by $H$, a
simple $Y$ is fixed on the right, and a simple $Z$ is fixed on both
sides. If $N=N_{XY}^Z>0$, then
\[
 \frac{\gcd(n,N)}{n}\FPdim(Z)
 \quad\text{is an algebraic integer.}
\]
\end{proposition}
\begin{proof}
As in Proposition~\ref{prop:cyclic-module-divisibility}, the fixed
simple trivializes the associator class on $H$, and
$H^2(C_n,\C^\times)=0$ supplies module structures. Thus
$A=\bigoplus_{h\in H}h$ is a connected separable algebra and
$X\otimes Y$ is an $A$-bimodule.
The $Z$-isotypic summand is stable under both actions, because $H$
fixes $Z$ on both sides. It is therefore an $A$-bimodule with underlying
object $NZ$. Its Frobenius--Perron dimension in the Morita dual is
$N\FPdim(Z)/n$, which is an algebraic integer.
Write $g=\gcd(n,N)=uN+vn$ with $u,v\in\Z$. Then
\[
 \frac{g}{n}\FPdim(Z)
 =u\frac{N}{n}\FPdim(Z)+v\FPdim(Z)
\]
is also an algebraic integer.
\end{proof}
This extends the prime-order argument without a pivotality assumption.
An application records a cyclic subgroup, its left and right actions,
the coefficient $N$, and an exact nonintegrality witness for the
scaled Frobenius--Perron dimension. The proof is independent of any
finite search or numerical approximation.
The extended checker finds $54$ witnesses among the $178$ tensors
surviving the basic passes. Of these, $49$ already satisfy the
prime-order test. The five additional order-four applications are:
\begin{center}\small
\begin{tabular}{@{}l r l r l@{}}
\toprule
Census ID & $g$ & $X=Y=Z$ & $N$ & Minimal polynomial of $\FPdim(Z)$\\
\midrule
$R7\_5\_4702$ & $b_1$ & $b_4$ & $1$ & $t^2-6t-4$\\
$R7\_5\_4704$ & $b_1$ & $b_4$ & $3$ & $t^2-8t-4$\\
$R7\_5\_4706$ & $b_1$ & $b_4$ & $5$ & $t^2-10t-4$\\
$R7\_5\_4757$ & $b_1$ & $b_4$ & $3$ & $t-10$\\
$R7\_5\_5676$ & $b_1$ & $b_3$ & $5$ & $t^2-10t-4$\\
\bottomrule
\end{tabular}
\end{center}
Each $N$ is odd, so the necessary quantity is $\FPdim(Z)/4$.
For a quadratic $t^2-at-4$, its minimal polynomial after division
by four is $t^2-(a/4)t-1/4$, which is not integral.
In the linear case the quotient is $5/2$.
The exact replay also verifies the subgroup actions and the minimal
polynomial of the Perron eigenvalue, using exact root isolation.
The application counts in this subsection may overlap earlier tests;
Table~\ref{tab:stages} charges each excluded ring only once.

\subsection{A cubic Fourier-module obstruction}
\begin{proposition}\label{prop:cubic-Fourier-module}
The fusion ring $R7\_5\_5007$ has no complex fusion-category categorification.
\end{proposition}
\begin{proof}
Write $g=b_1$, so $g^2=b_2$, and write the four remaining,
self-dual basis objects as $X_3,X_4,X_5,X_6$. Each is fixed by $g$.

The pointed subcategory generated by $g$ has trivial associator class:
the action on any one of these fixed simple objects supplies a rank-one
module category. Thus $A=1\oplus g\oplus g^2$ is a separable algebra.
The category of right $A$-modules has one simple over the free orbit
$\{1,g,g^2\}$ and three simples over each $X_i$. The dual group
$\widehat{C_3}$ acts by twisting the module structure, cyclically permuting
each triple and fixing the free-orbit simple. This action commutes with
the left action of the original category. For a primitive cube root
$\zeta$, its $\zeta$-eigenspace therefore yields a four-dimensional
representation $b_i\mapsto T_i$ of the fusion ring. In a Fourier basis,
\[
 (T_i)_{ab}=n_0+n_1\zeta+n_2\zeta^2,
 \qquad n_j\in\mathbb Z_{\ge0},\quad
 n_0+n_1+n_2=N_{i,a}^{b}, \qquad a,b\in\{3,4,5,6\}.
\]
Frobenius reciprocity for the module action gives $T_{i^*}=T_i^*$.
In particular the four matrices $T_i$ for $i=3,4,5,6$ are Hermitian.
The matrix $T_g$ is diagonal, with cube roots of unity on the diagonal.
Since $gX_i=X_i$, the identity $T_gT_i=T_i$ implies that a nontrivial
diagonal entry of $T_g$ forces the corresponding row of every $T_i$
($i=3,4,5,6$) to vanish. But a sum of three nonnegative multiples of
$1,\zeta,\zeta^2$ vanishes only when their total is divisible by three.
The entries $N_{3,3}^{4}=1$, $N_{3,4}^{3}=1$,
$N_{3,5}^{3}=2$, and $N_{3,6}^{3}=2$ consequently force $T_g=I$.

Set $T=T_3$. The following identity holds in the fusion ring, as checked
by multiplying the given integer matrices:
\[
 (X_3-9)(X_3+1)(X_3^3-X_3^2-4X_3-1)=6-3g-3g^2.
\]
Thus $(T-9I)(T+I)(T^3-T^2-4T-I)=0$.
The matrix $B=(N_{3,a}^{b})_{a,b=3}^6$ is
\[
 B=\begin{pmatrix}3&1&2&2\\1&2&2&3\\2&2&1&3\\2&3&3&3\end{pmatrix}.
\]
Its entries dominate the absolute values of the entries of $T$.
The leading principal minors of $9I-B$ are $6,41,268,47$, so
$\rho(B)<9$ and $9$ is not an eigenvalue of $T$.
The characteristic polynomial of a Hermitian matrix over
$\mathbb Q(\zeta)$ has rational coefficients; here they are algebraic
integers and hence integers. Since $x^3-x^2-4x-1$ is irreducible over
$\mathbb Q$, the only possibilities are $(x+1)^4$ and
$(x+1)(x^3-x^2-4x-1)$. Reduction modulo $(1-\zeta)$ shows
$\operatorname{tr}T\equiv\operatorname{tr}B=9\equiv0\pmod3$,
which excludes the first possibility. Consequently
\[
 \operatorname{tr}T=0,\qquad \operatorname{tr}T^2=10,
 \qquad \operatorname{tr}T^3=15.
\]
A real diagonal entry with total multiplicity $N$ has the form $N-3b$
with $0\le b\le\lfloor N/2\rfloor$. Together with the first two
moments, this forces diagonal $(0,-1,1,0)$. For an off-diagonal entry
with multiplicity $1$ or $2$, its squared norm is at least one; for
multiplicity three its squared norm is at least zero. There are four
of the former entries and two of the latter. Since the diagonal
squared norm is two, the second-moment bound is saturated: the four
entries have norm one, and the other two vanish. The only nonzero
triangle is on labels $3,4,5$. The diagonal contributions and the
edge contributions to $\operatorname{tr}T^3$ cancel, leaving six times
the real part of a product of three norm-one entries. Its absolute
value is at most six, contradicting $\operatorname{tr}T^3=15$.
The certificate checker also exhausts all 81 possible matrices after
these bounds, obtaining cubic traces only $-3$ and $6$.

\end{proof}


\subsection{Subring descent and involution modules}
A full fusion-closed, dual-stable subset of the basis generates a
fusion subcategory in any categorification. Thus the general rank-six
certificates from \cite{Rank6} give the following five ambient exclusions.
Each replay checks the restricted tensor and its nested exclusion
certificate, rather than relying on an ID or a recorded status alone.
\begin{center}\small
\begin{tabular}{@{}l l l@{}}
\toprule
Rank-seven ring & Rank-six subring & Ambient labels\\
\midrule
$R7\_5\_1033$ & $R6\_4\_304$ & $0,1,2,3,5,6$\\
$R7\_5\_1283$ & $R6\_3\_194$ & $0,1,2,3,5,6$\\
$R7\_5\_2019$ & $R6\_3\_244$ & $0,1,2,3,5,6$\\
$R7\_5\_5232$ & $R6\_2\_51$ & $0,1,2,4,5,6$\\
$R7\_5\_5681$ & $R6\_4\_19$ & $0,1,2,3,4,6$\\
\bottomrule
\end{tabular}
\end{center}

The anti-invariant representation of
Proposition~\ref{prop:anti-module} also applies to the extension.
Exact spectral, parity, and matrix bounds exclude
\[
 R7\_5\_1396,\ R7\_5\_1781,\ R7\_5\_2181,\
 R7\_5\_2247,\ R7\_5\_2454,\ R7\_5\_3078,\ R7\_5\_3102.
\]
The checker allows more than two fixed simples; the initial cheaper
involution pass was limited to smaller blocks. These are further
applications of the same general module-category argument.

The one-spectrum criterion of \cite[Theorem 6.8]{TPE} was also checked
on the full new band. It supplies $66$ exact witnesses, all on tensors
already excluded by the basic passes, so it contributes no additional
first-stage exclusions. The supplementary evidence includes its
independent replay and positive-model controls.

\subsection{Recorded outcome and open frontier}
Among the $5805$ new tensors, $5789$ are excluded, $1$ has a verified unitary model, and $15$ remain undecided.
The unresolved identifiers are
\begin{center}\small
\begin{tabular}{@{}l l l@{}}
\texttt{R7\_5\_443} & \texttt{R7\_5\_1214} & \texttt{R7\_5\_1964}\\
\texttt{R7\_5\_3895} & \texttt{R7\_5\_4530} & \texttt{R7\_5\_4531}\\
\texttt{R7\_5\_4569} & \texttt{R7\_5\_4653} & \texttt{R7\_5\_4982}\\
\texttt{R7\_5\_5029} & \texttt{R7\_5\_5039} & \texttt{R7\_5\_5040}\\
\texttt{R7\_5\_5343} & \texttt{R7\_5\_5623} & \texttt{R7\_5\_5778}\\
\end{tabular}
\end{center}
Each has its complete tensor and attempted checks in the companion data.
A feasible necessary system is not a categorical construction, and a
time-limited or incomplete search is not a nonexistence certificate.
No new tensor category is established for this frontier. The model in
Proposition~\ref{prop:m5-frobenius42} is the only new realized based-ring
type added to the catalog in this continuation.

